\documentclass[UTF8,fontset=fandol,zihao=5]{ctexart}

\usepackage[a4paper,margin=2cm]{geometry}
\usepackage{amsmath,amssymb,mathtools}
\usepackage{upgreek}
\usepackage[nolimits]{cmupint}
\usepackage{booktabs,tabularx,array}
\usepackage{xcolor}
\usepackage[most]{tcolorbox}
\usepackage{enumitem}
\usepackage{indentfirst}
\usepackage{hyperref}
\usepackage{float}

\usepackage{tikz}
\usepackage{diagbox}

\newcolumntype{W}[1]{%
  >{\hsize=#1\hsize\linewidth=\hsize
    \centering\arraybackslash}X}

\usetikzlibrary{
  positioning,
  calc,
  decorations.pathreplacing,
  arrows.meta
}

% relationbox 中知识关系图的统一样式
\tikzset{
  relnode/.style={
    font=\small,
    align=center
  },
  relline/.style={
    thick
  },
  relarrow/.style={
    -{Stealth},
    thick
  }
}

\DeclareTextFontCommand{\textbf}{\heiti\color{main}}

\definecolor{main}{HTML}{4C7FB3}
\definecolor{mainlight}{HTML}{F0F6FF}
\definecolor{method}{HTML}{8099B3}
\definecolor{methodlight}{HTML}{f0f6ff}

\hypersetup{colorlinks=true,linkcolor=main,urlcolor=main}
\setlength{\parindent}{2em}
\linespread{1.5}
\setlength{\abovedisplayskip}{5pt}
\setlength{\belowdisplayskip}{5pt}
\setlength{\abovedisplayshortskip}{3pt}
\setlength{\belowdisplayshortskip}{3pt}
\setlist[itemize]{leftmargin=2em,itemsep=0.1em,topsep=0.2em}
\setlist[enumerate]{leftmargin=2.3em,itemsep=0.1em,topsep=0.2em}

\newfontfamily\headingfont{TeX Gyre Heros}

\ctexset{
  section={format=\centering\heiti\headingfont\zihao{3}\color{main},
            name={第,章},number=\chinese{section},aftername=\quad,
            beforeskip=0.9em,afterskip=0.5em},
  subsection={format=\heiti\headingfont\zihao{5}\color{main},
              name={,、},number=\chinese{subsection},aftername=,
              beforeskip=0.6em,afterskip=0.22em},
  subsubsection={format=\heiti\headingfont\zihao{5}\color{main},
                  number={（\chinese{subsubsection}）},aftername=,
                  beforeskip=0.55em,afterskip=0.15em}
}

\newtcolorbox{relationbox}{
  enhanced,breakable,colback=mainlight,colframe=main,
  boxrule=0.65pt,arc=1.4mm,left=2mm,right=2mm,top=1.2mm,bottom=1.2mm
}
\newtcolorbox{methodbox}[1][]{
  enhanced,breakable,colback=methodlight,colframe=method,
  boxrule=0.65pt,arc=1.4mm,left=2mm,right=2mm,top=1.2mm,bottom=1.2mm,
  title=#1,fonttitle=\kaishu,fontupper=\kaishu,
  before upper={\setlength{\parindent}{2em}}
}

\newcommand{\jj}{\,\mathrm{j}}
\newcommand{\ee}{\,\mathrm{e}}
\newcommand{\CC}{\mathbb{C}}
\newcommand{\RR}{\mathbb{R}}
\newcommand{\QQ}{\mathbb{Q}}
\newcommand{\ZZ}{\mathbb{Z}}
\newcommand{\dd}{\,\mathrm{d}}
\newcommand{\ppi}{\,\uppi\,}
\newcommand{\FF}{\mathcal{F}}
\newcommand{\LL}{\mathcal{L}}
\newcommand{\Sa}{\operatorname{Sa}}
\newcommand{\RRe}{\operatorname{Re}}
\newcommand{\IIm}{\operatorname{Im}}
\newcommand{\dtft}{\operatorname{DTFT}}
\newcommand{\zf}{\mathcal{Z}}

\newfontfamily\russianfont{DejaVu Sans}

\begin{document}

\begin{center}
  {\heiti\headingfont\zihao{1}\color{main} 数字信号与系统笔记}
\end{center}

\section{模数转换}

\subsection{取样}

\begin{enumerate}
  \item 时域分析
  
  设取样间隔为$T_{\mathrm{s}}$，则取样信号为：

  \[
  \hat{x}(t)=x(t)\delta_T(t)=\sum_{k=-\infty}^{+\infty} x(kT_{\mathrm{s}})\delta_T(t-kT_{\mathrm{s}})
  \]

  \item 频域分析
  
  设取样角频率为$\varOmega_{\mathrm{s}}=\dfrac{2\uppi}{T_{\mathrm{s}}}$，
  由乘法的Fourier变换性质可知，取样信号的频域表示为：

  \[
  \hat{X}(\varOmega)=\dfrac{1}{T_{\mathrm{s}}}X(\varOmega)\ast P_T(\varOmega)=\dfrac{1}{T_{\mathrm{s}}}\sum_{k=-\infty}^{+\infty} X(\varOmega-k\varOmega_{\mathrm{s}})
  \]

  \item Nyquist采样定理
  
  取样后频谱副本之间不重叠的条件为：

  \[
  f_{\mathrm{s}} \leq 2f_{\max}
  \]

\end{enumerate}

\begin{methodbox}[【题型】求运算后信号的取样频率]
  
  将时域信号转换到频域后再判断$\varOmega_{\max}$的变化，如
  
  \begin{itemize}
    \item $f(at) \longrightarrow \dfrac{1}{|a|}F(\dfrac{\varOmega}{a})$，$a\varOmega_{\max}$
    \item $f_1(t) \ast f_2(t) \longrightarrow F_1(\varOmega)F_2(\varOmega)$，$\max\left\{\varOmega_{1\max},\,\varOmega_{2\max}\right\}$
    \item $f_1(t)f_2(t) \longrightarrow \dfrac{1}{2\uppi} F_1(\varOmega) \ast F_2(\varOmega)$，$\varOmega_{1\max} + \varOmega_{2\max}$
  \end{itemize}
  
\end{methodbox}

\begin{methodbox}
  注意：$\hat{x}(t)$是离散时间信号，仍然在时域$t-x(t)$上；而$x[n]$才是数字信号，在$n-x[n]$上。
\end{methodbox}

\subsection{内插}

\begin{enumerate}
  
  \item 内插公式
  
  \[
  x(t) = T\hat{x}(t) \ast 2f_{\mathrm{c}}\Sa(2\uppi f_{\mathrm{c}} t) = 2f_{\mathrm{c}}T\sum_{k=-\infty}^{+\infty} x(kT_{\mathrm{s}}) \Sa \left [ 2\uppi f_{\mathrm{c}}(t-kT_{\mathrm{s}}) \right ]
  \]

  \item 内插频率
  
  \[
  f_{\max} \leq f_{\mathrm{s}} \leq f_{\mathrm{s}}-f_{\max}
  \]

\end{enumerate}

\section{信号分析}

\subsection{时域分析}

\subsubsection{经典信号}

离散时间信号序列的角频率称为数字角频率，记为$\omega$，单位为rad，与角频率有关系$\omega=\varOmega T_{\mathrm{s}}$。

信号的常见表示法包括指针法、函数法、图像法等。

\begin{enumerate}

  \item 单位取样序列
  
  \[
  \delta[n]=
    \begin{cases}
    1, & n=0,\\
    0, & n\neq 0.
    \end{cases}
  \]

  \item 单位阶跃序列
  
  \[
  u[n]=
    \begin{cases}
    1, & n\geq 0,\\
    0, & n<0.
    \end{cases}
  \]

  \item 单位矩形序列
  
  \[
  R_N[n]=u[n]-u[n-N]=
    \begin{cases}
    1, & 0 \leq n \leq N-1,\\
    0, & n<0\;\text{或}\;n\geq N.
    \end{cases}
  =\sum_{k=0}^{N-1} \delta[n-k]
  \]

  \item 单位斜变序列
  
  \[
  r[n]=nu[n]=
    \begin{cases}
    n, & n\geq 0,\\
    0, & n<0.
    \end{cases}
  \]

  \item 实指数序列
  
  \[
  x[n]=a^n u[n], \quad a \in \mathbb{R}\; \text{且} \; a \ne 0
  \]
  
  \item 复指数序列
  
  \[
  x[n]=A \ee^{\jj (\omega_0 n + \varphi)}=A\cos(\omega_0 n + \varphi) + \jj A\sin(\omega_0 n + \varphi)
  \]
  
  \item 正弦序列
  
  \[
  x[n] = cos(\omega n) \text{或} x[n] = sin (\omega n)
  \]

  \begin{methodbox}[离散时间信号序列的周期性]
    若$x[n+N]=x[n]$，则称$x[n]$为周期序列，周期为$N$。

    若$x[n]$为正弦信号，则$x[n]$为周期信号当且仅当$\dfrac{2\uppi}{\omega}\in \QQ$，即$\omega = \dfrac{N}{m}\uppi, N,\,m \in \ZZ $。

    周期计算方式为$N=m\dfrac{2\uppi}{\omega},\,m \in \ZZ$。
  \end{methodbox}

\end{enumerate}

\subsubsection{运算}

\begin{enumerate}

  \item 几何变换
    
  几何变换可以用通式$\displaystyle y[n]=x\left [ a\left ( n - \dfrac{b}{a} \right ) \right ]$表示。
    
    \begin{enumerate}[label=(\arabic*)]
      \item 移位：$y[n]=x[n-k]$
      \item 反褶：$y[n]=x[-n]$
      \item 尺度变换：$y[n]=x[an]$
    \end{enumerate}

  \item 序列累加 $\displaystyle y[n]=\sum_{k=-\infty}^{n} x[k]$
  
  \item 复合运算
    \begin{enumerate}[label=(\arabic*)]
      \item 加法：$y[n]=x_1[n] + x_2[n]$
      \item 乘法：$y[n]=x_1[n] \cdot x_2[n]$
      \item 卷积：$y[n]=x[n] * h[n] = \displaystyle \sum_{k=-\infty}^{\infty} x[k] h[n-k]$
    \end{enumerate}
  
  \begin{methodbox}[【技巧】卷积计算]
    \begin{enumerate}[label=\arabic*.]
      \item 解析式法
      
      常用公式：
      
      \[
      \sum_{n=0}^{N-1} a^n = \frac{1-a^N}{1-a}, \qquad
      \sum_{n=0}^{\infty} a^n = \frac{1}{1-a}, \; |a|<1
      \]

      \[
      \sum_{n=0}^{N-1} n a^n = \frac{(N-1)a^{N+1}-Na^N+a}{(1-a)^2}, \qquad
      \sum_{n=0}^{\infty} n a^n = \frac{a}{(1-a)^2}, \; |a|<1
      \]

      \[
      \sum_{n=0}^{N-1} n = \frac{1}{2}N(N-1),\qquad
      \sum_{n=0}^{N-1} n^2 = \frac{1}{6}N(N-1)(2N-1)
      \]

      \item 图像法：反褶移位相乘求和，与连续时间卷积类似，不赘述。

      \item 排序法：将两列信号尾部对齐，对应位样值逐个相乘，按位相加，即得到结果。
    \end{enumerate}
  \end{methodbox}
  
\end{enumerate}

\subsection{频域分析}

\subsubsection{离散时间Fourier变换，DTFT}

\[
X(e^{\jj\omega}) = \FF \left\{ x[n] \right\} = \sum_{n=-\infty}^{\infty} x[n] e^{-\jj\omega n}
\]

特殊地，冲激信号$\delta[n]$的离散时间Fourier变换为$X(e^{\jj\omega})=1$。

其反变换为

\[
x[n] = \FF^{-1} \left\{ X(e^{\jj\omega}) \right\} = \int_{-\uppi}^{\uppi} X(e^{\jj\omega}) \ee^{\jj n\omega}\dd\omega
\]

\subsubsection{离散时间Fourier变换性质}

\begin{enumerate}
  \item 线性
  
  \[
  \FF \left\{ ax_1[n] + bx_2[n] \right\} = a \FF \left\{ x_1[n] \right\} + b \FF \left\{ x_2[n] \right\} 
  \]     

设$\displaystyle \FF \{ x[n] \} = X(\omega)$，则有

  \item 共轭对称性
  
  \[
  \FF \left\{ x^\ast(t) \right\} = X^\ast(-\omega)
  \]

  当$x(t)$为实信号时，有$\displaystyle X(-\varOmega) = X^\ast(\varOmega)$。此时，$x(t)$额外有奇偶虚实性、

    \begin{center}
    \renewcommand{\arraystretch}{1.5}
    \begin{tabular}{|c|c|}
    \hline
    实偶：$\RRe X=\RRe X^\ast$
    &
    虚奇：$\IIm X=-\IIm X^\ast$
    \\
    \hline
    幅偶：$|X|=|X^\ast|$
    &
    相奇：$\varphi(X)=-\varphi(X^\ast)$
    \\
    \hline
    \end{tabular}
    \end{center}

  \item 运算性质
    \begin{center}
    \begin{tabularx}{0.98\textwidth}{>{\bfseries}p{1.5cm}X X}
    \toprule
    运算 & Fourier变换 & Fourier逆变换 \\
    \midrule

    平移
    &
    $\displaystyle
    \FF \{ x[n-n_0] \}
    =
    X(\ee^{\jj\omega})\ee^{-\jj n_0 \omega}
    $
    &
    $\displaystyle
    \FF^{-1} \{ X(\ee^{\jj(\omega-\omega_0)}) \}
    =
    x[n]\ee^{\jj \omega_0 n}
    $
    \\

    卷积
    &
    $\displaystyle
    \FF \{ x[n]\ast y[n] \}
    =
    X(\ee^{\jj\omega})Y(\ee^{\jj\omega})
    $
    &
    $\displaystyle
    \FF^{-1}
    \{ X(\ee^{\jj\omega})\ast Y(\ee^{\jj\omega}) \}
    =x[n]y[n]
    $
    \\

    \bottomrule
    \end{tabularx}
    \end{center}

  \item Parseval定理
  
  \[
  \sum_{-\infty}^{+\infty} |x[n]|^2 = \frac{1}{2\uppi} \int_{-\uppi}^{\uppi} |X(\ee^{\jj\omega})|^2 \dd \omega
  \]

\end{enumerate}

\subsection{复频域分析}

\subsubsection{Z变换}

令$z=\ee^{\sigma + \jj \omega}=r\ee^{\jj \omega}$，

\begin{enumerate}
  \item 定义
\[
\zf\left\{x[n]\right\}=\sum_{n=-\infty}^{\infty} x[n] z^n
\]
\[
\zf^{-1}\left\{X(z)\right\}=\oint\frac{1}{2\uppi\jj} X(z) z^{n-1} \dd z
\]

  \item 收敛域（Region of Convergence，ROC）
  
  使Z变换收敛的$z$的取值范围称为收敛域。

  不同信号的收敛域如下：
  \begin{enumerate}[label=(\arabic*)]
    \item 右边信号：$|z|>R_1$
    \item 左边信号：$|z|<R_2$
    \item 双边信号：$R_1<|z|<R_2$
    \item 时限信号：$0<|z|<\infty$
  \end{enumerate}
\end{enumerate}

\subsubsection{基本信号的Z变换}

\begin{center}
\begin{table}[H]
  \centering
  \begin{tabularx}{\textwidth}{>{\bfseries}p{3cm} c c}
    \toprule
    序列 & 时域序列 & Z变换 \\
    \midrule
    单位冲激序列 & $x[n]=\delta[n]$
      & $X(z)=1,\quad 0<|z|<\infty$ \\[6pt]

    单位右边阶跃序列 & $x[n]=u[n]$
      & $X(z)=\dfrac{1}{1-z^{-1}},\quad |z|>1$ \\[6pt]

    单位左边阶跃序列 & $x[n]=-u[-n-1]$
      & $X(z)=\dfrac{1}{1-z^{-1}},\quad 0<|z|<1$ \\[6pt]

    斜坡信号 & $x[n]=nu[n]$
      & $X(z)=\dfrac{z}{(1-z^{-1})^2},\quad |z|>1$ \\[6pt]

    右边余弦信号 & $x(t)=\cos(\varOmega_0t)u(t)$
      & $X(s)=\dfrac{s}{s^2+\varOmega_0^2},
        \quad \operatorname{Re}(s)>0$ \\

    右边正弦信号 & $x(t)=\sin(\varOmega_0t)u(t)$
      & $X(s)=\dfrac{\varOmega_0}{s^2+\varOmega_0^2},
        \quad \operatorname{Re}(s)>0$ \\
    \bottomrule
  \end{tabularx}
\end{table}
\end{center}

\subsubsection{Z变换性质}

\begin{table}[H]
  \centering
  \begin{tabularx}{\textwidth}{
    >{\bfseries}p{2cm}
    >{\centering\arraybackslash}X
    >{\centering\arraybackslash}p{3.2cm}
  }
    \toprule
    性质 & 变化 & ROC 变化 \\
    \midrule
    线性
    & $\zf\left\{a x_1[n]+b x_2[n]\right\}=aX_1(z)+bX_2(z)$
    & $\text{ROC} \supseteq R_1 \cap R_2$ \\

    反褶
    & $\zf\left\{x[-n]\right\}=X\left(\dfrac{1}{z}\right)$
    & 不变 \\

    尺度变换
    & $\zf\left\{x[a^n n]\right\}=X\left(\dfrac{z}{a}\right)$
    & 缩放为 $|a||z|$ \\

    时移
    & $\zf\left\{x[n-n_0]\right\}=X(z)z^{-n_0}$
    & 不变 \\

    共轭
    & $\zf\left\{x^*[n]\right\}=X^*(z^*)$
    & 不变 \\

    复频域微分
    & $\zf\left\{nx*[n]\right\}=-z\dfrac{\dd X(z)}{\dd z}$
    & 不变 \\

    时域卷积
    & $\zf\left\{x_1[n]*x_2[n]\right\}=X_1(z)X_2(z)$
    & $\text{ROC} \supseteq R_1 \cap R_2$ \\

    复频域卷积
    & $\displaystyle\zf\left\{x_1[n]x_2[n]\right\}=\dfrac{1}{2\uppi\jj} \oint X_1(v)X_2\left(\dfrac{z}{v}\right)\dd v$
    &  \\
    \bottomrule
  \end{tabularx}
\end{table}

\begin{enumerate}
  \item 初值定理 $\displaystyle x[0]=\lim_{z\to\infty}X(z)$
  \item 终值定理 $\displaystyle\lim_{n\to\infty}x[n]=\lim_{z\to 1}[(z-1)X(z)]$
\end{enumerate}

\begin{methodbox}[【题型】求Z反变换]
  \begin{itemize}
    \item 幂级数展开法（长除法）
    \item 部分分式展开法
    
    对于多项式信号$X(s)=\dfrac{P(s)}{Q(s)}$，首先将其化为整式与真分式，
    随后分别处理。
    \begin{enumerate}[label=\arabic*.]
      \item 整式：$X_0(s)=\sum_{i=0}^{n}c_i s^i \longleftrightarrow x_0(t)=\sum_{i=0}^{n}c_i \delta^{(i)}(t)$
      \item 真分式：将真分式$X_1(s)$展开为部分分式$\displaystyle\sum_{i=1}^{n}\frac{K_i}{(s-p_i)^r}$，分别作Laplace反变换
      
      \[
      \LL^{-1}\left\{\frac{K_i}{(s-p_i)^r}\right\} = K_i \frac{t^{r-1}}{(r-1)!} e^{p_i t} u(t),\, s.t.\, \RRe s > \RRe p_i
      \]
      \[
      \LL^{-1}\left\{\frac{K_i}{(s-p_i)^r}\right\} = K_i \frac{t^{r-1}}{(r-1)!} e^{p_i t} [-u(-t)],\, s.t.\, \RRe s < \RRe p_i
      \]
      其中$K_i$可由留数法求得。
      
      \[
      K_i=\left.X_1(s)(s-p_i)^r\right|_{s=p_i}
      \]
      
      注意到左边信号只需要把右边信号的$u(t)$换为$-u(-t)$即可，故后续只以右边信号为例。

      其中，特殊极点类型有特殊的方法来计算Laplace反变换。下列例中均以右边信号为例：
      \begin{itemize}
        \item $p_i$为单实数极点，分式形如$\dfrac{K_i}{s-p_i}$
        
        \[
        \LL^{-1}\left\{X(s)\right\}=K_i u(t) e^{p_i t}
        \]
        
        \item $p_i$为重极点，分式形如$\displaystyle\sum_{k=1}^{r}\dfrac{K_{i,k}}{(s-p_i)^k}$：
        
        \[
        \LL^{-1}\left\{X(s)\right\}=K_{i,k} \dfrac{t^{k-1}}{(k-1)!} u(t) e^{p_i t}
        \]
        
        \item $p_{1,\,2}=\alpha\pm\varOmega_0$为共轭复数极点，$K=C+\jj D$，分式形如$\dfrac{K}{s-p_1}+\dfrac{K^\ast}{s-p_2}$

        \[
        \LL^{-1}\left\{X(s)\right\}=2e^{\alpha t} \left[C\cos (\varOmega_0 t) - D\sin (\varOmega_0 t)\right]  u(t) 
        \]

      \end{itemize}
    \end{enumerate}
    \item 留数法
  \end{itemize}
\end{methodbox}

\subsubsection{Z变换和离散Fourier变换、Laplace变换关系}

\section{系统分析}

系统分析主要围绕冲激响应$h(t)$、$H(\jj \varOmega)$和$H(s)$进行。

\subsection{系统}

\subsubsection{表示}

\begin{itemize}
  \item 基本运算单元
  \item 线性常系数微分方程
\end{itemize}

\subsubsection{基本概念}
系统是对输入信号进行加工、变换、处理，并产生输出信号的物理装置。

\subsubsection{冲激响应与阶跃响应}

冲激响应是系统对单位冲激信号的响应，记为$h(t)$。冲激响应是系统的特征函数，能够完全描述系统的性质。

阶跃响应是系统对单位阶跃信号的响应，记为$g(t)$。

\begin{methodbox}[二者关系]
  \[
  g(t)=\int_{-\infty}^{t}h(\tau)\dd\tau=h(t) \ast u(t),\, h(t)=\frac{\dd g(t)}{\dd t}
  \]
\end{methodbox}

任意输入信号$x(t)$的输出信号$y(t)$可以表示为冲激响应的卷积，即
\[
y(t) = x(t) \ast h(t)
\]

\subsubsection{典型冲激响应}

\begin{center}
\begin{tabularx}{0.98\textwidth}{>{\bfseries}p{2cm}X}
\toprule
系统类型 & 冲激响应 $h(t)$ \\
\midrule
        标量乘法器 & $h(t)=a\delta(t)$ \\
        微分器 & $h(t)=\dfrac{\mathrm{d}\delta(t)}{\mathrm{d}t}=\delta'(t)$ \\
        积分器 & $h(t)=\displaystyle\int_{-\infty}^{t}\delta(\tau)\,\mathrm{d}\tau=u(t)$ \\
        延时器 & $h(t)=\delta(t-\tau)$ \\
\bottomrule
\end{tabularx}
\end{center}

\subsubsection{分类}

\begin{enumerate}
  
  \item 线性系统与非线性系统
  
  线性系统满足齐次性和叠加性，亦即映射$\mathcal{T} : x(t) \mapsto y(t)$是线性算子。
  
  \begin{methodbox}[线性系统的判定]
  
    \[\mathcal{T}\left\{ax_1(t) + bx_2(t)\right\} = a\mathcal{T}\left\{x_1(t)\right\} + b\mathcal{T}\left\{x_2(t)\right\}\]
  
  一般线性系统的响应可以写成各输入信号的线性组合。

  加法、微分、积分等运算都是线性算子。
  
  注意：表达式带常数的系统一般不是线性系统，例如$y(t)=2x(t)+3$。
  可用$\mathcal{T}\left\{x(0)\right\}=0$快速判断。
  
  \end{methodbox}

  \item 时变系统与时不变系统
  
  时不变系统的运行规则不会随时间推移发生改变，亦即映射$\mathcal{T} : x(t) \mapsto y(t)$不是$t$的函数。

  \begin{methodbox}[时不变系统的判定]
    \[\mathcal{T}\left\{x(t-t_0)\right\} = y(t-t_0)\]
    
    常见反例：$y(t)=x(2t)$，$y(t)=tx(t)$。
  \end{methodbox}

  \item 因果系统与非因果系统
  
  因果系统的输出仅由当前和过去的输入决定，亦即$y(t_0)$只依赖于$x(t), \, t \in \left ( -\infty, t_0 \right ]$。

  \begin{methodbox}[因果系统的判定]
    因果系统冲激响应满足$h(t)=0,\, t<0$的一般都是非因果系统。
    
    一般包含$x(t+a),\, a>0$的都是非因果系统，如$y(t)=x(t+1)$。
  \end{methodbox}

  \item 稳定系统与不稳定系统

  稳定系统（BIBO稳定）在接入有限大小的输入信号时，能够产生有限范围的输出信号。

  \begin{methodbox}[稳定系统的判定]
    \[\int_{-\infty}^{+\infty} h(t) \dd t < \infty\]
    
    常见反例：$h(t)=sin(\omega_0 t)u(t)$。
  \end{methodbox}

  \begin{methodbox}[因果稳定系统的判定]
    零极点图上的极点均位于左半平面。
  \end{methodbox}

\end{enumerate}

\subsection{频域与复频域分析}
\begin{center}
\begin{tabularx}{0.98\textwidth}{>{\bfseries}p{2cm}X X X}
\toprule
分析方式 & 特性 & 微分方程形式 & 零极点形式 \\
\midrule
频域分析 & 
$H(\jj \varOmega) = \dfrac{Y_\mathrm{ZS}(\varOmega)}{X(\varOmega)}$ & 
$H(\jj \varOmega) = \dfrac{\displaystyle\sum_{j=0}^{m}b_j\left(\jj \varOmega\right)^j}{\displaystyle\sum_{i=0}^{n}a_i\left(\jj \varOmega\right)^i}$ &
$H(\jj \varOmega)=H_0\frac{\displaystyle\prod_{j=1}^{m}(\jj \varOmega-z_j)}{\displaystyle\prod_{i=1}^{n}(\jj \varOmega-p_i)}$ \\
复频域分析 & 
$H(s) = \dfrac{Y_\mathrm{ZS}(s)}{X(s)}$ & 
$H(s) = \dfrac{\displaystyle\sum_{j=0}^{m}b_j s^j}{\displaystyle\sum_{i=0}^{n}a_i s^i}$ &
$H(s)=H_0\frac{\displaystyle\prod_{j=1}^{m}(s-z_j)}{\displaystyle\prod_{i=1}^{n}(s-p_i)}$ \\
\bottomrule
\end{tabularx}
\end{center}

\begin{methodbox}[【注意】]
  不要忘记零极点形式中的$H_0$！
\end{methodbox}

\begin{methodbox}[频域分析的其他零极点形式]
  零点矢量可以写作$\jj\varOmega-z_j=N_j e^{\jj\psi_j}$，极点矢量可以写作$\jj\varOmega-p_i=M_i e^{\jj\theta_i}$。
  
  由两矢量的指数形式可得
  
  \[
  H(\jj\varOmega) = \left|H(\jj\varOmega)\right|\ee^{-\jj\varOmega\varphi(\varOmega)} = H_0\frac{\displaystyle\prod_{j=1}^{m}N_j}{\displaystyle\prod_{i=1}^{n}M_i}\ee^{\displaystyle\jj\left(\displaystyle\sum_{k=1}^{m}\psi_k-\displaystyle\sum_{l=1}^{n}\theta_{l}\right)}
  \]

  其中$\left|H(\jj\varOmega)\right|= H_0\frac{\displaystyle\prod_{j=1}^{m}N_j}{\displaystyle\prod_{i=1}^{n}M_i}$称为幅频响应，
  $\displaystyle\varphi(\varOmega)=\displaystyle\sum_{k=1}^{m}\psi_k-\displaystyle\sum_{l=1}^{n}\theta_{l}$称为相频响应。
\end{methodbox}

\begin{methodbox}[由零极点图求幅频响应与相频响应]
  由零极点图可得
  \[
  \left|H(\jj\varOmega)\right|= H_0\frac{\displaystyle\prod_{j=1}^{m}N_j}{\displaystyle\prod_{i=1}^{n}M_i}
  ,\qquad
  \varphi(\varOmega)=\displaystyle\sum_{k=1}^{m}\psi_k-\displaystyle\sum_{l=1}^{n}\theta_{l}
  \]

  其中$N_j$为零点矢量长度，$\psi_j$为零点矢量与实轴正向的夹角，$M_i$为极点矢量长度，$\theta_i$为极点矢量与实轴正向的夹角。
  
\end{methodbox}

\begin{methodbox}[【技巧】从微分方程推得频域冲激响应]
  做变换$\displaystyle \frac{\dd^n}{\dd t^n} \longrightarrow (\jj \varOmega)^n$，即可得到频域冲激响应。例：

  \[
  \frac{\dd^2}{\dd t^2}y(t)+6\frac{\dd}{\dd t}y(t)+5y(t)=\frac{\dd}{\dd t}x(t)+2x(t)
  \]
  \[
  \left[(\jj \varOmega)^2+6(\jj \varOmega)+5\right]Y(\varOmega)=(\jj \varOmega+2)X(\varOmega)
  \]
  \[
  H(\jj\varOmega)=\frac{Y(\varOmega)}{X(\varOmega)}=\frac{(\jj \varOmega)^2+6(\jj \varOmega)+5}{\jj \varOmega+2}
  \]

\end{methodbox}

\subsection{信号流图}

\subsubsection{从信号流图求冲激响应}

原理过于复杂不叙述，这里只介绍一般的做题方法。

从$X(s)$出发，将能够到达$Y(s)$的最简路径全部列出，路径内部相乘路径之间相加（后称写成线性形式）作为分子；
从$X(s)$出发，将能够到达$X(s)$的最简路径（包括1）全部列出，写成线性形式作为分母；
这样便得到了冲激响应$H(s)$。

此外还有Mason公式

\[
H(s)=\frac{1}{\varDelta}\sum_{i=1}^{n}{P_i\varDelta_i}
\]

不多赘述。

\subsubsection{从冲激响应求信号流图}

\begin{enumerate}
  \item 直接法
  \[
  H(s)=H_1(s) \cdot H_2(s)
  \]
  \begin{methodbox}[用直接法做信号流图]
    两类直接法的作图过程为求冲激响应的逆过程，具体不赘述。
    \begin{itemize}
      \item 直接法\uppercase\expandafter{\romannumeral1}：先画$b$，再画$a$。
      
      \item 直接法\uppercase\expandafter{\romannumeral2}：先画$a$，再画$b$。
    \end{itemize}
  \end{methodbox}
  \item 级联法
  \[
  H(s)=\prod_{i=1}^{n}H_i(s)
  \]
  优点：方便调控零点与极点

  缺点：误差累积较大
  
  \item 并联法
  \[
  H(s)=\sum_{i=1}^{n}H_i(s)
  \]
  
  优点：方便调控极点，同时减小误差累积。

  缺点：不方便调控零点。
\end{enumerate}

\section{滤波器设计}

\subsection{无失真传输}

无失真传输要求幅频响应为常数，且相频响应为线性函数。

\begin{enumerate}
  \item 时域条件：
  \[
  h(t)=K\delta(t-t_0)
  \]
  \item 频域条件：
  \[
  \left|H(\jj\varOmega)\right|=K, \qquad \varphi(\varOmega)=\varOmega t_0
  \]
\end{enumerate}

\subsection{理想模拟低通滤波器}

\begin{enumerate}
  \item 响应函数特征：

  \begin{itemize}
    \item 时域响应：$h(t)=K\Sa(t-t_0)$
    \item 幅频响应：$\left|H(\jj\varOmega)\right|=K\left(u(\varOmega_c)-u(-\varOmega_c)\right)$
    \item 相频响应：$\varphi(\varOmega)=\varOmega t_0$
  \end{itemize}

\begin{methodbox}[【理解】理想模拟低通滤波器不存在]
  理想模拟低通滤波器是非因果系统（冲激响应为取样函数），现实中无法制造。
\end{methodbox}

  \item 系统函数特征

  \begin{methodbox}[【题型】由平方响应函数求系统函数]
    \begin{enumerate}[label=(\arabic*)]
      \item 作代换$s=\jj\varOmega$
      \item 求零点与极点
      \item 取左半平面极点写出$H(s)$
      
      例：$\left|H(\jj\varOmega)\right|^2=\dfrac{4(\varOmega^4+64)}{\varOmega^4-10\varOmega^2+169}$
      
      解：
      \[
      H(s)H(-s)=\left.\frac{4(\varOmega^4+64)}{\varOmega^4-10\varOmega^2+169}\right|_{\varOmega=-\jj s}=\frac{4(s^4+64)}{s^4+10s^2+169}
      \]
      
      解得零点$s_\text{z}=\pm2\pm\jj 2$，极点$s_\text{p}=\pm2\pm\jj 2$。
      选取左半平面极点$s_\text{p}=-2\pm\jj 2$，则
      \[
      H(s)=H_0\frac{s^2+4s+8}{s^2+4s+13}
      \]

      又根据$\left.H(s)\right|_{s=0}=\left.H(\jj\varOmega)\right|_{\varOmega=0}$，
      得$H_0=2$，则

      \[
      H(s)=2\frac{s^2+4s+8}{s^2+4s+13}
      \]
    \end{enumerate}
  \end{methodbox}
\end{enumerate}
\subsection{现实模拟低通滤波器}

\subsubsection{Butterworth滤波器}

Butterworth逼近
\[
\left|H(\jj\varOmega)\right|^2=\dfrac{1}{1+\left(\dfrac{\varOmega}{\varOmega_c}\right)^{2N}}
\]

\begin{methodbox}[设计Butterworth滤波器]
  
  需要参数：$A_p$、$\varOmega_p$、
  $A_s$、$\varOmega_s$、$\varOmega_c$
  
  有如下关系：
  \[
  A_p=10\lg\left(\frac{\varOmega_p}{\varOmega_c}\right)
  \]
  \[
  A_s=10\lg\left(\frac{\varOmega_s}{\varOmega_c}\right)
  \]

  \begin{enumerate}
    \item 确定滤波器阶数$N$
    \[
    N=\left\lceil\frac{\lg{\dfrac{10^{0.1A_p}-1}{10^{0.1A_s}-1}}}{\lg {\dfrac{\varOmega_p}{\varOmega_s}}}\right\rceil
    \]

    \item 求滤波器极点
    \[
    p_k=\exp\left[\displaystyle\jj\dfrac{\uppi}{2}+\displaystyle\jj\dfrac{(2k+1)\uppi}{2N}\right], \quad k=0,\,1,\,\ldots,\,2N-1
    \]

    \item 求滤波器传递函数
    
    取左半平面极点，即$k=0,\,1,\,\ldots,\,N-1$
    \[
    H(s)=\frac{\varOmega_c^N}{\displaystyle\prod_{k=0}^{N-1}(s-p_k\varOmega_c)}
    \]

    也可查表求传递函数分母，常见分母如下
    \begin{itemize}
      \item $N$=1：$p+1$
      \item $N$=2：$p^2+\sqrt{2}p+1$
      \item $N$=3：$(p^2+p+1)(p+1)$
    \end{itemize}
  \end{enumerate}
\end{methodbox}

\subsubsection{\russianfont Чебышёв滤波器（略）}

\subsubsection{椭圆滤波器（略）}

\subsection{模拟频率变换}

\begin{table}[htbp]
    \centering
    \label{tab:filter-transform}
    \renewcommand{\arraystretch}{2}
    \setlength{\tabcolsep}{6pt}
    \begin{tabular}{@{}clc@{}}
        \toprule
        滤波器类型 &
        \shortstack{归一化低通滤波器 $H_{\mathrm{LP}}(p)$} &
        \shortstack{需要变换 $p=q(s)$} \\
        \midrule
        \shortstack{低通$\leftrightarrow$低通\\
        $H_{\mathrm{LP}}(s)\rightarrow H_{\mathrm{LP}}(p)$} &
        $\begin{aligned}
            \lambda_p &= \frac{1}{a},\quad
            \lambda_s = \frac{\Omega_s}{a\Omega_p}
        \end{aligned}$ &
        $\displaystyle p=\frac{s}{a\Omega_p}$ \\
        \addlinespace
        \shortstack{低通$\leftrightarrow$高通\\
        $H_{\mathrm{HP}}(s)\rightarrow H_{\mathrm{LP}}(p)$} &
        $\begin{aligned}
            \lambda_p &= \frac{1}{a},\quad
            \lambda_s = \frac{\Omega_p}{a\Omega_s}
        \end{aligned}$ &
        $\displaystyle p=\frac{\Omega_p}{as}$ \\
        \addlinespace
        \shortstack{低通$\leftrightarrow$带通\\
        $H_{\mathrm{LP}}(p)\rightarrow H_{\mathrm{BP}}(s)$} &
        $\begin{aligned}
            \lambda_p &= \frac{1}{a},\quad
            \lambda_s = \frac{\Omega_{s2}-\Omega_{s1}}
            {a(\Omega_{p2}-\Omega_{p1})}
        \end{aligned}$ &
        $\displaystyle p=\frac{s^2+\Omega_0^2}{aBs}$ \\
        \addlinespace
        \shortstack{低通$\leftrightarrow$带阻\\
        $H_{\mathrm{LP}}(p)\rightarrow H_{\mathrm{BS}}(s)$} &
        $\begin{aligned}
            \lambda_p &= \frac{1}{a},\quad
            \lambda_s = \frac{\Omega_{p2}-\Omega_{p1}}
            {a(\Omega_{s2}-\Omega_{s1})}
        \end{aligned}$ &
        $\displaystyle p=\frac{Bs}{a(s^2+\Omega_0^2)}$ \\
        \bottomrule
    \end{tabular}
\end{table}

\end{document}
